\documentclass[UTF8,a4paper,11pt,fontset=windows]{ctexart}
\usepackage[margin=2.1cm]{geometry}
\usepackage{graphicx,float,booktabs,tabularx,fancyhdr,xurl,fvextra}
\usepackage[hidelinks]{hyperref}
\setlength{\parskip}{.25em}
\setlength{\emergencystretch}{2em}
\renewcommand{\arraystretch}{1.2}
\pagestyle{fancy}\fancyhf{}
\fancyhead[L]{互联网数据库开发 · 知行学习团队}
\fancyhead[R]{2413581 黄奕浩}\fancyfoot[C]{\thepage}
\setlength{\headheight}{14pt}\raggedbottom
\begin{document}
\begin{titlepage}\thispagestyle{empty}\centering
\vspace*{.6cm}\includegraphics[width=.43\textwidth]{../../figures/nankai_logo_color.png}
\vspace{.9cm}\rule{.88\textwidth}{.8pt}\par\vspace{.55cm}
{\fontsize{27pt}{35pt}\selectfont\bfseries 互联网数据库开发课程设计\par}
\vspace{.65cm}{\fontsize{21pt}{29pt}\selectfont\bfseries 需求文档\par}
\vspace{.55cm}\rule{.88\textwidth}{.8pt}\par
\vspace{1.3cm}{\Large 知行学习团队\par}
\vspace{.6cm}{\large 黄奕浩\quad 2413581\quad 组长\par
王乐之\quad 2411302\par}
\end{titlepage}
\section{项目目标与用户}
知行学习团队门户与资源协作平台用于介绍团队和成员、分享学习内容、组织活动并汇集课程成果。访客可以认识团队、浏览已发布内容和下载作业；登录成员可以维护自己的资料，并使用资源收藏、评论和活动报名；管理员负责内容发布、资料维护和审核。

系统以前台内容展示与后台管理组成。团队及个人介绍、留言交流、信息编辑与信息发布是课程设计的核心内容。数据库设计为15张业务表，采用Yii2 Advanced组织MVC代码，MySQL保存数据。团队按模块分工，每人负责至少一个Model、Controller和View。

\section{功能需求}
\subsection{公共账户与团队展示}
注册使用用户名、邮箱和密码，用户名不可重复。普通注册账户固定为member角色，不接受客户端指定管理员身份。登录后可查看当前身份并退出；后台单独校验管理员权限。密码保存为哈希，页面不展示密码和验证密钥。

首页展示团队介绍、成员卡片与最新公告。成员列表提供姓名、学号和个人介绍，详情页展示完整内容。资料编辑使用会话中的用户ID，普通成员只修改自己的记录；管理员可以在后台维护所有成员以及团队介绍。无效成员编号返回404。

\subsection{公告发布与作业下载}
管理员可以新增、编辑、发布、撤下和删除公告。公告分为草稿、已发布和已撤下状态，只有已发布且发布时间不晚于当前时间的公告可被访客读取。删除操作使用POST并校验CSRF令牌；用户输入经转义后显示。

下载中心按团队材料与个人作业分类。记录保存相对于data目录的文件路径和SHA256摘要，下载前验证路径范围、文件是否存在及摘要是否一致。不可用文件不显示有效下载按钮。团队材料位于data/team，个人材料位于data/personal，SQL位于data/install.sql。

\subsection{资源、活动与交流}
资源、动态图表、活动报名和资源评论由王乐之负责；团队留言审核回复、公共操作日志和邮件服务由本人负责。团队留言、审核、日志和邮件已实现；资源、活动与资源评论保留待接入状态。

\section{数据需求与接口}
公共模块涉及user、team、member\_profile、announcement、assignment\_file。资源模块涉及resource\_category、resource、tag、resource\_tag、resource\_favorite。活动交流涉及activity、activity\_registration、guestbook\_message、resource\_comment、audit\_log。

公共身份接口site/me只返回登录状态、用户ID、用户名与角色，不返回密码哈希、邮箱或auth\_key。资源发布者、评论作者和报名用户均使用会话身份，不能相信表单提交的用户编号。成员与账户一对一，资源标签和收藏用唯一约束去重，活动报名用活动与用户组合唯一键去重。

\section{页面与操作约束}
前台导航包含团队首页、成员、公告和作业下载。后台围绕团队介绍、成员维护和公告管理提供操作入口。表单显示字段标签、校验错误和保存反馈；移动端在窄屏下按单列排列成员卡片。

所有写请求执行CSRF校验，后台动作在服务器端检查管理员身份。查询通过ActiveRecord或绑定参数执行，正文使用文本转义。注册使用的邮箱只作为账号信息，本项目不依赖外部邮件服务，当前注册流程不发送邮件。

\section{分工与交付}
本人负责公共账户、团队成员、公告下载、留言审核回复、公共邮件日志、集成部署及需求、实现、部署文档和用户手册，并负责完整录屏的组织、录制、讲解、剪辑与提交。王乐之负责资源、活动、资源评论、动态图表、数据库设计、设计文档和最终PPT统稿，提供本人模块演示素材并配合录制。PPT与录屏均作为重点交付预留充分时间，按实际代码、文档和演示产物说明工作。

团队最终提交需求、设计、实现、用户手册、部署、项目展示和5至15分钟录屏。Git用于版本控制，Tower用于任务管理。每个PHP文件写明用途与责任归属，文件命名统一，示例公告用于演示发布流程。各人保留本模块代码、必要截图及提交记录。

\end{document}
